\documentclass[11pt,letterpaper]{article}
\usepackage[T1]{fontenc}
\usepackage[utf8]{inputenc}
\usepackage{lmodern}
\usepackage[margin=1in,headheight=15pt]{geometry}
\usepackage{amsmath,amssymb,amsthm,mathtools,mathrsfs}
\usepackage{microtype,booktabs,longtable,array,enumitem}
\usepackage[dvipsnames]{xcolor}
\usepackage{fancyhdr}
\usepackage[colorlinks=true,linkcolor=MidnightBlue,citecolor=MidnightBlue,urlcolor=MidnightBlue]{hyperref}
\usepackage[nameinlink,noabbrev]{cleveref}
\hypersetup{pdftitle={Finite Hahn Deformations in Surcomplex Analysis},pdfauthor={Research draft prepared with OpenAI ChatGPT},pdfsubject={Multivariable division, conservation of multiplicity, and residue duality}}
\pagestyle{fancy}\fancyhf{}
\fancyhead[L]{\small Finite Hahn deformations}
\fancyhead[R]{\small Surcomplex analysis}
\fancyfoot[C]{\thepage}
\setlength{\parskip}{0.3em}\setlength{\parindent}{1.2em}
\setlist{nosep,leftmargin=2em}
\setcounter{tocdepth}{2}
\newtheorem{theorem}{Theorem}[section]
\newtheorem{lemma}[theorem]{Lemma}
\newtheorem{proposition}[theorem]{Proposition}
\newtheorem{corollary}[theorem]{Corollary}
\theoremstyle{definition}
\newtheorem{definition}[theorem]{Definition}
\newtheorem{example}[theorem]{Example}
\newtheorem{hypothesis}[theorem]{Hypothesis}
\crefname{hypothesis}{Hypothesis}{Hypotheses}
\Crefname{hypothesis}{Hypothesis}{Hypotheses}
\theoremstyle{remark}
\newtheorem{remark}[theorem]{Remark}
\crefname{theorem}{Theorem}{Theorems}
\crefname{lemma}{Lemma}{Lemmas}
\crefname{proposition}{Proposition}{Propositions}
\crefname{corollary}{Corollary}{Corollaries}
\crefname{remark}{Remark}{Remarks}
\newcommand{\No}{\mathbf{No}}
\newcommand{\SC}{\mathbf{SC}}
\newcommand{\C}{\mathbb C}
\newcommand{\R}{\mathbb R}
\newcommand{\N}{\mathbb N}
\newcommand{\Q}{\mathbb Q}
\newcommand{\OO}{\mathcal O}
\newcommand{\HH}{\mathcal H}
\newcommand{\mm}{\mathfrak m}
\newcommand{\NF}{\mathcal N_F}
\newcommand{\RF}{\mathcal R_F}
\newcommand{\supp}{\operatorname{supp}}
\newcommand{\st}{\operatorname{st}}
\newcommand{\Tr}{\operatorname{Tr}}
\newcommand{\Spec}{\operatorname{Spec}}
\newcommand{\Hom}{\operatorname{Hom}}
\newcommand{\Res}{\operatorname{Res}}
\newcommand{\rank}{\operatorname{rank}}
\newcommand{\Id}{\operatorname{id}}
\newcommand{\im}{\operatorname{im}}
\newcommand{\length}{\operatorname{length}}
\newcommand{\sh}{^{\sharp}}
\newcommand{\HInt}{\mathop{\int^{\!H}}\nolimits}
\newcommand{\vH}{v_H}
\newcommand{\eps}{\varepsilon}
\newcommand{\dif}{\mathrm d}
\newcommand{\qvec}{\mathbf q}
\newcolumntype{L}[1]{>{\raggedright\arraybackslash}p{#1}}
\numberwithin{equation}{section}
\begin{document}
\begin{titlepage}
\centering
\vspace*{0.7cm}
{\Large\scshape A multivariable continuation of the supplied framework\par}
\vspace{0.9cm}
{\Huge\bfseries Finite Hahn Deformations\par}
\vspace{0.3cm}
{\LARGE\bfseries in Surcomplex Analysis\par}
\vspace{0.6cm}
{\Large Multivariable Division, Conservation of Multiplicity,\par and Residue Duality\par}
\vspace{0.8cm}
{\large September 21, 2026\par}
\vspace{0.6cm}
\begin{minipage}{0.92\textwidth}
\begin{abstract}
The supplied manuscript proposes a several-variable theory of Hahn-coherent analytic sets with support control on a fixed ordinary domain. We establish its isolated-complete-intersection case. If an ordinary holomorphic system $f_1,\ldots,f_n$ on a polydisc has an isolated common zero of multiplicity $\mu$, then every strictly positive Hahn perturbation $F_i=f_i+E_i$ admits division with a unique $\mu$-dimensional remainder space, without repeatedly shrinking the coefficient domain. All output supports lie in the input support plus the additive monoid generated by the perturbation supports. The quotient over the Hahn valuation ring is finite free of rank $\mu$, even for higher-rank value groups and non-Noetherian valuation rings.

We identify its geometric points with the actual surcomplex zeros in the ordinary zero's monad and identify its local factors with formal local algebras at the displaced zeros. Consequently their multiplicities sum exactly to $\mu$. Coefficientwise Grothendieck contours give a perfect residue pairing, a trace--Jacobian identity, a moving-simple-zero residue formula, and a discriminant identity. A valuation estimate controls changes of normal forms, residues, and spectral invariants under changes of the perturbation. Two coupled two-variable examples, one polynomial and one transcendental, exhibit finite total residues assembled from individually infinite terms.

The proof uses classical Koszul exactness, the homological perturbation lemma, and ordinary residue duality. The contribution offered here is the explicit fixed-domain Hahn implementation and its geometric and quantitative consequences, not a claim to have invented those classical mechanisms. Full proofs of the extension steps are given; the finite examples are also checked symbolically.
\end{abstract}
\end{minipage}
\vfill
\begin{minipage}{0.92\textwidth}\small\raggedright
\textbf{Research status.} This is an original research draft developed from the attached \texttt{surcomplex\_analysis(3).tex}. The main combined statements below were not located in the literature checked for this draft; that is a qualified novelty assessment, not a priority certification. The work answers a precise unfinished direction of the supplied manuscript, not an identified longstanding published conjecture. It is not peer reviewed or proof-assistant verified.
\end{minipage}
\end{titlepage}
\tableofcontents
\newpage

\section{The question being answered}\label{sec:intro}

\subsection{From one variable to a finite analytic space}
The supplied text \cite{Companion} defines Hahn-coherent functions by a well-ordered family of ordinary holomorphic coefficient functions on one common ordinary domain. It proves one-variable preparation and counts the zeros that split from a multiple ordinary zero. Its final research directions ask for multivariable coherent division, analytic-set theory, and finite-mapping theorems with support control that is not concealed in successive local choices.

The central question of this article is the following.
\begin{quote}
If an isolated ordinary holomorphic system of $n$ equations in $n$ variables is perturbed at arbitrarily many well-ordered infinitesimal scales, does its entire finite local algebra survive, on one fixed coefficient domain, with its multiplicities and residue pairing intact?
\end{quote}
We prove that it does. The hypothesis is not that the ordinary zero is simple. Arbitrary isolated complete intersections are allowed, including nonreduced ones. Nor is the perturbation required to be polynomial, to have finitely many scales, or to belong to a discrete or rank-one valuation.

The resulting theorem is stronger than the assertion that some nearby roots exist. It constructs a finite free algebra containing all of them, determines the sum of their formal local multiplicities, transports every Koszul relation, and gives a residue functional that recovers normal-form coefficients and multiplication matrices.

\subsection{The main package in informal notation}
Let $D\subset\C^n$ be a polydisc centered at $0$, and suppose that
\[
 f_1=\cdots=f_n=0\quad\hbox{in }D
 \qquad\Longleftrightarrow\qquad z=0.
\]
Write $\mu$ for the ordinary local intersection multiplicity. Work in a set-sized Hahn field
\[
 K=\C((t^\Gamma)),\qquad t^\gamma=\omega^{-\gamma},
\]
where $\Gamma$ is a divisible ordered subgroup of $\No$. Its valuation ring is $R$, and its infinitesimal ideal is $\mm$. Put
\[
 F_i=f_i+E_i,\qquad E_i=\sum_{\gamma>0} e_{i,\gamma}(z)t^\gamma,
 \quad e_{i,\gamma}\in\OO(D),
\]
with well-ordered supports. Let $E$ be their finite union and $M=E^*$ its finite-sum monoid, including $0$.

Choose polynomial representatives $q_1=1,q_2,\ldots,q_\mu$ of a basis of the ordinary quotient. Then every Hahn-holomorphic $G$ has
\begin{equation}\label{eq:main-division-intro}
 G=\sum_{a=1}^{\mu}c_aq_a+\sum_{i=1}^n F_iQ_i,
 \qquad
 \supp(c_a),\ \supp(Q_i)\subseteq\supp(G)+M.
\end{equation}
The remainder is unique; the individual quotients need not be. All coefficients of every $Q_i$ are holomorphic on the original $D$.

The zeros $\xi$ are finite in number, all infinitesimal, and
\begin{equation}\label{eq:main-count-intro}
 \sum_{F\sh(\xi)=0}
 \dim_K\frac{K[[u_1,\ldots,u_n]]}
 {\bigl(\widehat F_{1,\xi},\ldots,\widehat F_{n,\xi}\bigr)}
 =\mu.
\end{equation}
Here $\widehat F_{i,\xi}$ denotes the Taylor expansion at the \emph{actual displaced zero} $\xi$, not at its ordinary standard part.

The coefficientwise residue functional $\RF$ has the identities
\begin{equation}\label{eq:main-residue-intro}
 \RF(GJ_F)=\Tr(M_G),\qquad
 \RF(J_F)=\mu,
 \qquad
 J_F=\det\left(\frac{\partial F_i}{\partial z_j}\right).
\end{equation}
If the zeros are simple, then
\begin{equation}\label{eq:main-simple-intro}
 \RF(G)=\sum_{F\sh(\xi)=0}\frac{G\sh(\xi)}{J_F\sh(\xi)}.
\end{equation}
These formulas remain meaningful when individual summands on the right are infinitely large and cancel to a finite result.

\subsection{What is inherited, and what is extended}
The terminology \emph{Hahn-coherent}, its coefficientwise contour interpretation, and the distinction between strong summation and fine-topological convergence come from the supplied manuscript. We use its conventions, but all several-variable constructions needed below are stated explicitly. No unproved several-variable assertion from that manuscript is assumed.

There is substantial prior analysis over surreal fields: Alling's monograph already treats formal series in finitely many surreal and surcomplex variables \cite{Alling}. General analytic structures on valued fields, including Weierstrass-type machinery, are developed by Cluckers and Lipshitz \cite{CluckersLipshitz}. The strong-summation algebra comes from Hahn--Mal'cev--Neumann theory \cite{Neumann,Poonen}. The perturbation formulas used below are classical homological perturbation formulas \cite{Crainic}, and finite local algebras and residues have a mature ordinary theory \cite{Demailly,Kunz,CattaniDickenstein}.

Accordingly, none of the following is advertised as new in isolation: the implicit function theorem, the Koszul complex, the perturbation lemma, ordinary conservation of number, or ordinary Grothendieck duality. The proposed additional result is their precise realization in the full common-domain Hahn coefficient algebra, together with the comparison to actual surcomplex points and the explicit support and precision estimates. A search under a new phrase is not enough to certify novelty; \cref{sec:status} records the more limited conclusion justified by the literature check.

\section{Hahn-coherent functions in several variables}\label{sec:hahn}

\subsection{The set-sized field and its valuation}
Fix an ordinary integer $n\ge1$. All algebra in the proofs takes place in sets. Fix a divisible ordered subgroup $\Gamma\subseteq(\No,+)$ and set
\[
 K=\left\{\sum_{\gamma\in S}c_\gamma t^\gamma:
 c_\gamma\in\C,\ S\subseteq\Gamma\text{ well ordered}\right\}.
\]
The convention $t^\gamma=\omega^{-\gamma}$ embeds $K$ into $\SC=\No[i]$. Every specified set of surcomplex coefficients and supports can be accommodated by enlarging $\Gamma$ to a set-sized divisible subgroup. Algebraic closedness of $K$ follows from the divisible-value-group, algebraically-closed-residue-field theorem for Hahn fields \cite[Corollary 4]{Poonen}.

For nonzero $a\in K$, let $v(a)=\min\supp(a)$, and put $v(0)=+\infty$. Define
\[
 R=\{a:v(a)\ge0\},\qquad
 \mm=\{a:v(a)>0\}.
\]
The residue map $R\to\C$ takes the coefficient at exponent $0$. In the surcomplex embedding, $R$ consists of the finite elements of $K$, and $\mm$ of its infinitesimals. We use no rank-one norm and no Noetherian property of $R$.

\subsection{Strong summation and coefficient modules}
\begin{definition}[Strongly summable family]
A set-indexed family of Hahn series is strongly summable if the union of its supports is well ordered and each exponent receives contributions from only finitely many members of the family. Its sum is formed coefficient by coefficient.
\end{definition}
The same definition applies to series with coefficients in an arbitrary complex vector space $W$. We write $W((t^\Gamma))$ for that space of series. A complex-linear map $L:W\to W'$ extends coefficientwise and satisfies
\begin{equation}\label{eq:linear-support}
 \supp(Lw)\subseteq\supp(w).
\end{equation}
It also preserves all strongly summable families. Neither continuity nor a bound on the sizes of complex coefficients is needed for this assertion.

\begin{lemma}[Support calculus]\label[lemma]{lem:neumann}
Let $E\subseteq\Gamma_{>0}$ be well ordered and let
\[
 E^*=\{0\}\cup\{e_1+\cdots+e_k:k\ge1,\ e_i\in E\}.
\]
Then $E^*$ is well ordered, and a fixed exponent has only finitely many representations as a finite word in $E$. If $S\subseteq\Gamma$ is well ordered, then $S+E^*$ is well ordered, and a fixed exponent has only finitely many representations $s+e_1+\cdots+e_k$ with $s\in S$.
\end{lemma}
This is the positive-support form of Neumann's lemma \cite{Neumann}; the additional set $S$ uses the usual finite-decomposition rule for the sum of two well-ordered subsets of an ordered abelian group. Permutations of a given finite word cause only finitely many repetitions.

\begin{remark}[What the lemma does not say]\label[remark]{rem:rank}
The intersection $E^*\cap(-\infty,\eta)$ need not be finite. For example, in a group containing $1$ and $\omega$, all ordinary positive integers are less than $\omega$. Thus $1+t+t^2+\cdots$ is a valid Hahn sum, although its ordinary finite partial sums do not approximate it to valuation $>\omega$. Throughout this article, infinite formulas are justified by strong summation, never by silently treating $k\,\min E$ as cofinal in $\Gamma$.
\end{remark}

\subsection{The analytic coefficient algebra}
For a polydisc $D\subset\C^n$, define
\begin{align*}
 A&=\OO(D),\\
 A_H&=A((t^\Gamma)),\\
 A_H^+&=\left\{G\in A_H:\supp(G)\subseteq\Gamma_{\ge0}\right\}.
\end{align*}
For $G\ne0$, $\vH(G)$ is the least exponent with a nonzero \emph{coefficient function}; thus $\vH(G)$ is not computed at a particular ordinary point. Define $\vH(0)=+\infty$. The ring $A_H^+$ is an $R$-algebra, and $A_H$ is a $K$-algebra. The latter is generally much larger than the algebraic tensor product $A\otimes_\C K$.

Let
\[
 D\sh_K=\{c+\eps:c\in D,\ \eps\in\mm^n\}.
\]
For $G=\sum_\gamma g_\gamma t^\gamma$, its coherent evaluation is
\begin{equation}\label{eq:evaluation}
 G\sh(c+\eps)=
 \sum_{\gamma}\sum_{\alpha\in\N^n}
 \frac{\partial^\alpha g_\gamma(c)}{\alpha!}
 t^\gamma\eps^\alpha.
\end{equation}
Zero coordinates of $\eps$ are harmless. The multi-index notation is
$\alpha!=\prod_j\alpha_j!$ and $\eps^\alpha=\prod_j\eps_j^{\alpha_j}$.

\begin{proposition}[Evaluation and differentiation]\label[proposition]{prop:evaluation}
The sum in \eqref{eq:evaluation} is strongly summable. Evaluation is an injective algebra homomorphism, and it commutes with every partial derivative. If $G\in A_H^+$, then
\[
 \st G\sh(c+\eps)=g_0(c).
\]
\end{proposition}
\begin{proof}
Let $P$ be the finite union of the positive supports of the nonzero $\eps_j$. Each term in \eqref{eq:evaluation} has support in $\supp(G)+P^*$. A fixed exponent has only finitely many decompositions into an exponent of $G$ and a word in $P$, by \cref{lem:neumann}; there are only finitely many assignments of that word to the $n$ coordinates. Hence the sum is strong. The product and derivative rules reduce coefficientwise to the ordinary multivariable Taylor rules; all regroupings have the same finite-contribution justification. At an ordinary point $c$, the value is $\sum_\gamma g_\gamma(c)t^\gamma$. If this vanishes for every $c$, each coefficient function vanishes. Finally all Taylor terms except $g_0(c)$ have positive valuation in the nonnegative-support case.
\end{proof}

\subsection{Divided differences at infinitesimal points}
The following observation will rule out extraneous algebraic characters of our finite quotient.
\begin{lemma}[Coherent Taylor division]\label[lemma]{lem:taylor-division}
Let $D$ be centered at $0$, let $\xi\in\mm^n$, and let $G\in A_H$. There exist $H_j\in A_H$ such that
\begin{equation}\label{eq:divided-difference}
 G(z)-G\sh(\xi)=\sum_{j=1}^n(z_j-\xi_j)H_j(z).
\end{equation}
More generally, for every ordinary integer $N\ge1$,
\begin{equation}\label{eq:finite-taylor}
 G(z)-\sum_{|\alpha|<N}
 \frac{\partial^\alpha G\sh(\xi)}{\alpha!}(z-\xi)^\alpha
 \ \in\ (z_1-\xi_1,\ldots,z_n-\xi_n)^N A_H.
\end{equation}
The coefficient functions in the remainders are holomorphic on $D$.
\end{lemma}
\begin{proof}
For an ordinary holomorphic $g$, coordinate telescoping writes
\[
 g(z)-g(w)=\sum_j(z_j-w_j)h_j(z,w),
\]
where each $h_j$ is holomorphic on $D\times D$. The apparent singularity in the one-coordinate divided difference is removable. Expanding $h_j$ in $w$ at $0$ and substituting $w=\xi$ produces Hahn coefficients holomorphic on the entire $z$-polydisc $D$. If $P$ is the union of the positive supports of $\xi_j$, their supports lie in $P^*$. Applying this to all $g_\gamma$ and multiplying by $t^\gamma$ gives \eqref{eq:divided-difference} with support in $\supp(G)+P^*$.

Finite Taylor's formula with holomorphic remainder on a convex polydisc is obtained by repeated divided differences, or by the ordinary segment-integral remainder on $D\times D$. Substitute $w=\xi$ in its coefficient functions in exactly the same fashion. Only finitely many degree-$N$ factors $(z-\xi)^\alpha$ are involved, and \cref{lem:neumann} justifies their coefficientwise evaluation. This proves \eqref{eq:finite-taylor}.
\end{proof}

\section{The ordinary complete intersection and its contraction}\label{sec:ordinary}

\begin{hypothesis}[Fixed ordinary system]\label[hypothesis]{hyp:ordinary}
The polydisc $D\subset\C^n$ is centered at $0$; $f_1,\ldots,f_n\in\OO(D)$; and their common zero set in $D$ is exactly $\{0\}$. Put
\[
 V=A/(f_1,\ldots,f_n),\qquad \mu=\dim_\C V.
\]
\end{hypothesis}
We first justify that this global quotient really is the desired finite ordinary local algebra.

\begin{proposition}[The ordinary quotient and Koszul resolution]\label[proposition]{prop:ordinary-koszul}
Under \cref{hyp:ordinary}, restriction to the germ at $0$ induces
\begin{equation}\label{eq:ordinary-local}
 V\simeq\OO_{\C^n,0}/(f_1,\ldots,f_n),\qquad 1\le\mu<\infty.
\end{equation}
The Koszul complex $C_\bullet=K_\bullet(f;A)$ is exact in positive homological degrees, and its degree-zero homology is $V$.
\end{proposition}
\begin{proof}
The local ring $\OO_{\C^n,0}$ is regular of dimension $n$. Since the common zero is isolated, the ideal of the $f_i$ is primary to the maximal ideal. Thus the $f_i$ form a system of parameters and hence a regular sequence. At any other point of $D$, one $f_i$ is a unit in the local ring, and its Koszul complex is contractible. Consequently the Koszul complex of coherent analytic sheaves is exact in positive degrees everywhere on $D$. Its degree-zero quotient sheaf is supported at $0$ and has the finite-dimensional stalk in \eqref{eq:ordinary-local}.

The polydisc is Stein. Cartan's theorem B, applied to the coherent kernels and images of this finite complex, makes the complex of global sections exact in the same degrees and makes the global quotient map onto that supported sheaf surjective. A sheaf supported at one point has its finite stalk as its global sections. This proves both assertions. These are ordinary complex-analytic facts; standard foundations and the Stein vanishing theorem can be found in \cite[Chapters II and IX]{Demailly}.
\end{proof}

Choose monomials $q_1=1,q_2,\ldots,q_\mu$ whose classes form a basis of $V$. Such a choice exists because some power of the ordinary maximal ideal is zero in $V$, so finitely many monomials span it. Define $i:V\to C_0=A$ by these representatives and $p:C_0\to V$ by quotienting; set $p=0$ on positive degrees.

The Koszul differential, in homological grading, is
\begin{equation}\label{eq:koszul-differential}
 d(a\,e_{j_1}\wedge\cdots\wedge e_{j_k})
 =\sum_{r=1}^k(-1)^{r-1}f_{j_r}a\,
 e_{j_1}\wedge\cdots\widehat e_{j_r}\cdots\wedge e_{j_k}.
\end{equation}
There is a complex-linear degree-$+1$ operator $h$ satisfying
\begin{equation}\label{eq:contraction}
 dh+hd=1-ip,\qquad pi=1,\qquad hi=0,\quad ph=0,\quad h^2=0.
\end{equation}

\begin{proof}[Construction of $h$]
Let $B_k=\im(d:C_{k+1}\to C_k)$. In degree $0$, take
$C_0=i(V)\oplus B_0$. In each positive degree, exactness allows a complex-vector-space complement $L_k$ to $B_k$ such that $d:L_k\to B_{k-1}$ is an isomorphism. Define $h$ on $B_{k-1}$ as its inverse and set $h=0$ on $L_k$ and on $i(V)$. The identities in \eqref{eq:contraction} follow on the displayed direct-sum components.
\end{proof}

\begin{remark}[Why an algebraic splitting is useful]\label[remark]{rem:splitting}
The map $h$ need not be $A$-linear or continuous in a Fr\'echet topology. It maps holomorphic functions on $D$ to holomorphic functions on the \emph{same} $D$, because it is chosen inside the complex of global sections. Extending $h$ coefficientwise preserves support by \eqref{eq:linear-support}. Thus no analytic norm estimate and no sequence of smaller polydiscs is needed to iterate it along positive Hahn exponents. This is the fixed-domain feature on which the proof depends.

The existence proof uses ordinary choices of complex-vector-space complements. It is not, by itself, an effective algorithm for an arbitrary black-box holomorphic system. In concrete systems, ordinary division operators or the residue reconstruction in \cref{sec:residue} supply usable normal forms.
\end{remark}

\section{A support-controlled perturbation lemma}\label{sec:perturbation}
Extend $C_\bullet,d,h,i,p$ coefficientwise to Hahn modules. Let
\[
 F_j=f_j+E_j,\qquad \supp(E_j)\subseteq\Gamma_{>0},
\]
and let $\delta$ be the Koszul differential built from the $E_j$. Then the perturbed differential is
\[
 \partial=d+\delta,
 \qquad \partial^2=0,
 \qquad \delta i=0.
\]
The last equality simply says that a homological differential vanishes on degree $0$.

Put $E=\bigcup_j\supp(E_j)$ and $M=E^*$. When all $E_j=0$, use $E=\varnothing$ and $M=\{0\}$.

\begin{lemma}[Hahn perturbation with degree-zero homology]\label[lemma]{lem:hpl}
Define
\begin{align}
 S&=(1+\delta h)^{-1}=\sum_{k\ge0}(-\delta h)^k,
 &\mathcal A&=S\delta,\label{eq:resolvent}\\
 p_F&=pS=p-p\mathcal A h,
 &h_F&=hS=h-h\mathcal A h.\label{eq:perturbed-maps}
\end{align}
Every displayed sum is strong. These maps satisfy
\begin{equation}\label{eq:perturbed-contraction}
 \partial h_F+h_F\partial=1-ip_F,
 \qquad p_F\partial=0,
 \qquad p_Fi=1.
\end{equation}
For an input supported in $S_0$, both $p_F$ and $h_F$ have output supported in $S_0+M$.
\end{lemma}
\begin{proof}
Each occurrence of $\delta$ adds an exponent from $E$, while $h,d,p,i$ do not add exponents. Thus the $k$th summand in \eqref{eq:resolvent} has support in $S_0+kE$, and \cref{lem:neumann} gives well-ordering and finite contributions over all $k$. The same argument applies to $\mathcal A$. In particular all algebraic rearrangements below are legitimate coefficientwise finite identities.

This is a specialized form of the classical homological perturbation lemma \cite{Crainic}. For clarity, we verify the necessary identities with our sign convention. The geometric inverses give
\begin{equation}\label{eq:A-identities}
 \delta h\mathcal A=\mathcal A h\delta=\delta-\mathcal A,
 \qquad
 \mathcal A=(1+\delta h)^{-1}\delta
 =\delta(1+h\delta)^{-1}.
\end{equation}
The relation $d\delta+\delta d+\delta^2=0$ implies
\[
 (1+\delta h)(d\mathcal A+\mathcal A d)(1+h\delta)
 =-\delta ip\delta.
\]
Indeed the left side expands to
\[
 d\delta+\delta d+\delta(hd+dh)\delta
 =-\delta^2+\delta(1-ip)\delta.
\]
Multiplying by the two inverses yields
\[
 d\mathcal A+\mathcal A d=-\mathcal A ip\mathcal A.
\]
Since $\delta i=0$, also $\mathcal A i=0$, so
\begin{equation}\label{eq:A-anticommute}
 d\mathcal A+\mathcal A d=0.
\end{equation}
Now \eqref{eq:A-identities} gives
\begin{align*}
 \partial h_F+h_F\partial
 &=1-ip+\mathcal A h+h\mathcal A
      -dh\mathcal A h-h\mathcal A hd\\
 &=1-ip+ip\mathcal A h
      +h\mathcal A(1-dh-hd)\\
 &=1-ip+ip\mathcal A h
 =1-ip_F.
\end{align*}
In the second line we used \eqref{eq:A-anticommute}; the third uses $\mathcal A i=0$. Similarly,
\begin{align*}
 p_F\partial
 &=p\delta-p\mathcal A hd-p\mathcal A h\delta\\
 &=p\mathcal A(1-hd)
 =p\mathcal A(ip+dh)=0,
\end{align*}
using $pd=0$ and \eqref{eq:A-anticommute}. Finally $p_Fi=pi-p\mathcal A hi=1$. The support estimate was established in the first paragraph.
\end{proof}

\begin{remark}[Not a new abstract perturbation lemma]
Invertibility of $1+\delta h$ is already the appropriate algebraic smallness condition in the classical lemma. What has been checked here is that this inverse exists on the full Hahn modules, respects one ordinary coefficient domain, and carries the exact support estimate required by the surcomplex application.
\end{remark}

\section{Fixed-domain division and finite freeness}\label{sec:division}

\begin{theorem}[Hahn-coherent complete-intersection division]\label[theorem]{thm:division}
Assume \cref{hyp:ordinary} and let $F=f+E$ be as above. For every $G\in A_H$ there are $Q_1,\ldots,Q_n\in A_H$ and a unique $c=(c_1,\ldots,c_\mu)\in K^\mu$ such that
\begin{equation}\label{eq:division}
 G=\sum_{a=1}^\mu c_aq_a+\sum_{j=1}^nF_jQ_j.
\end{equation}
One choice is
\begin{equation}\label{eq:normal-form}
 c=\NF(G)=p(1+\delta h)^{-1}G,
 \qquad (Q_1,\ldots,Q_n)=h(1+\delta h)^{-1}G.
\end{equation}
The chosen quotients and the remainder satisfy
\begin{equation}\label{eq:division-support}
 \supp(c_a),\ \supp(Q_j)\subseteq\supp(G)+M,
 \qquad
 v(c_a),\ \vH(Q_j)\ge\vH(G).
\end{equation}
If $G\in A_H^+$, then $c_a\in R$ and $Q_j\in A_H^+$. In particular
\begin{equation}\label{eq:finite-free}
 B_R:=A_H^+/(F_1,\ldots,F_n)\simeq R^\mu,
 \qquad
 B_K:=A_H/(F_1,\ldots,F_n)\simeq K^\mu
\end{equation}
as modules, with basis the classes of the $q_a$.
\end{theorem}
\begin{proof}
Apply \eqref{eq:perturbed-contraction} to a degree-zero input $G$. Since $\partial G=0$,
\[
 G=ip_FG+\partial h_FG,
\]
which is precisely \eqref{eq:division}. The support and valuation estimates follow from \cref{lem:hpl}, since $M\subseteq\Gamma_{\ge0}$.

If a remainder $i(c)$ belongs to the ideal of the $F_j$, it is a degree-zero boundary: $i(c)=\partial Q$. Applying $p_F$ and using \eqref{eq:perturbed-contraction} yields $c=0$. Thus the remainder is unique and the kernel of $p_F$ is exactly that ideal. The same identities restrict to nonnegative supports, which proves the assertion over $R$ as well as over $K$.
\end{proof}

\begin{corollary}[All Koszul relations lift]\label[corollary]{cor:koszul}
The Koszul complex $K_\bullet(F;A_H^+)$ has zero homology in positive degrees. The same holds over $A_H$. If a cycle $Z$ has support $S_0$, then the explicit preimage $h_FZ$ has support in $S_0+M$ and coefficients on $D$.
\end{corollary}
\begin{proof}
For a positive-degree cycle, $p_FZ=0$ by degree, so \eqref{eq:perturbed-contraction} reads $Z=\partial h_FZ$.
\end{proof}

\subsection{Algebra structure and reduction}
The module isomorphism in \eqref{eq:finite-free} transports multiplication to $R^\mu$. Its structure constants are
\begin{equation}\label{eq:structure}
 q_aq_b\equiv\sum_{c=1}^\mu C_{ab}^cq_c\pmod{(F)},
 \qquad C_{ab}^c=\NF(q_aq_b)_c.
\end{equation}
They have support in $M$, and their exponent-zero coefficients are the ordinary structure constants of $V$. Hence
\begin{equation}\label{eq:special-fiber}
 B_R/\mm B_R\simeq V.
\end{equation}
This is a genuine finite free deformation of the ordinary finite algebra. In the usual algebraic sense, $\Spec B_R\to\Spec R$ is finite flat of rank $\mu$. That statement is not a claim about proper maps in an undeclared topology on the class $\SC^n$.

\subsection{Dependence on the splitting and the domain}
Although the selected quotients in \eqref{eq:normal-form} depend on $h$, the remainder map does not depend on $h$ once the representative subspace $\operatorname{span}_\C\{q_a\}$ has been fixed: uniqueness in \cref{thm:division} proves this immediately. Different quotients differ by a Koszul syzygy, as \cref{cor:koszul} describes.

If $D'\subseteq D$ is a smaller centered polydisc on which the same ordinary system has only the zero $0$, the same representatives form a basis of its ordinary quotient. Restricting \eqref{eq:division} to $D'$ and using uniqueness shows that the remainder is unchanged. Thus shrinking once for a contour construction later does not change the finite algebra or its normal forms.

\subsection{A coefficient recursion with honest scope}
Let $W=(1+\delta h)^{-1}G$. Its coefficients obey
\begin{equation}\label{eq:recursion}
 W_\nu=G_\nu-\sum_{e\in E}\delta_e hW_{\nu-e},
 \qquad
 \NF(G)=pW,\quad Q=hW.
\end{equation}
Only finitely many displayed terms contribute to a fixed $\nu$ in the supported recursion. Iterating the formula yields the finite-word expansion of \cref{lem:hpl}. This is an exact coefficient algorithm when ordinary operations and support membership are effective. It is not a blanket claim that every valuation truncation can be listed in finite time; \cref{rem:rank} explains the obstruction.

\begin{example}[Why division at the ordinary center is wrong]
For $F(z)=z^2-t$, the formal series $F$ is a unit in $K[[z]]$, because its constant coefficient $-t$ is nonzero in the field $K$. Consequently $K[[z]]/(z^2-t)=0$. Nevertheless $F$ has two infinitesimal zeros $\pm t^{1/2}$. The correct cluster algebra is $\OO(D)((t^\Gamma))/(z^2-t)$, which has dimension $2$. Only after locating an actual zero do we pass to its formal local ring. The distinction is essential in \cref{eq:main-count-intro}.
\end{example}

\section{Actual surcomplex zeros and conservation of multiplicity}\label{sec:zeros}

\subsection{Multiplication matrices force infinitesimal coordinates}
For $G\in A_H$, let $M_G$ denote multiplication by its class on $B_K$, in the basis $q_1,\ldots,q_\mu$. Its $a$th column is $\NF(Gq_a)$. In particular the coordinate matrices $M_j=M_{z_j}$ belong to $\operatorname{Mat}_\mu(R)$ and reduce to multiplication by $z_j$ on $V$.

\begin{lemma}[Coordinate characteristic polynomials]\label[lemma]{lem:coordinate-polynomials}
Each
\[
 P_j(T)=\det(TI-M_j)\in R[T]
\]
is monic of degree $\mu$ and satisfies
\[
 P_j(T)\bmod\mm=T^\mu.
\]
Every root of $P_j$ in $K$, or in any valued field extension with the same ordered comparison, is infinitesimal.
\end{lemma}
\begin{proof}
The ordinary local algebra $V$ has nilpotent maximal ideal; multiplication by each $z_j$ is therefore nilpotent. This proves the reduction assertion. Write $P_j(T)=T^\mu+\sum_{k<\mu}a_kT^k$, with $a_k\in\mm$. If a nonzero root $\lambda$ had $v(\lambda)\le0$, division by $\lambda^\mu$ would yield
\[
 1+\sum_{k<\mu}a_k\lambda^{k-\mu}=0,
\]
where every term in the sum has positive valuation. This is impossible. The root $0$ is already infinitesimal.
\end{proof}

\begin{theorem}[The finite zero set is the spectrum]\label[theorem]{thm:points}
Evaluation induces a bijection between common zeros of $F\sh$ in $D\sh_K$ and $K$-algebra characters $B_K\to K$. Every such zero lies in $\mm^n$. There are at least one and at most $\mu$ distinct zeros, and there are no additional common zeros in the full surcomplex thickening $D\sh\subset\SC^n$ outside $K^n$.
\end{theorem}
\begin{proof}
If $z=c+\eps$ is a coherent zero, \cref{prop:evaluation} gives $f_i(c)=0$ for every $i$, so $c=0$. Evaluation at $z\in\mm^n$ is an algebra homomorphism that annihilates the $F_i$, hence a character of $B_K$.

Conversely, let $\chi:B_K\to K$ be a character and put $\xi_j=\chi(z_j)$. The Cayley--Hamilton identities $P_j(z_j)=0$ hold in $B_K$, so \cref{lem:coordinate-polynomials} gives $\xi\in\mm^n$. For every $G\in A_H$, apply \cref{lem:taylor-division} and then $\chi$ to obtain
\[
 \chi([G])=G\sh(\xi).
\]
Taking $G=F_i$ shows that $\xi$ is a zero. This also proves uniqueness of the character with a given coordinate tuple. A nonzero finite-dimensional commutative algebra over the algebraically closed field $K$ has finitely many maximal ideals, with residue fields $K$; their number lies between $1$ and its dimension $\mu$.

Finally, the identities $P_j(z_j)\in(F)$ are identities in $A_H$, so they can be evaluated at any surcomplex zero after enlarging the set-sized Hahn workspace to include that point. Each coordinate must then be a root of a polynomial over $K$. Since $K$ is algebraically closed, all its polynomial roots in any field extension already belong to $K$. Thus no extra zeros appear by passing to the full class field.
\end{proof}

\subsection{Comparison with the formal germs at the moved points}
For a zero $\xi$, define
\begin{equation}\label{eq:local-taylor}
 \widehat F_{i,\xi}(u)=
 \sum_{\alpha\in\N^n}
 \frac{\partial^\alpha F_i\sh(\xi)}{\alpha!}u^\alpha
 \in K[[u_1,\ldots,u_n]].
\end{equation}
This is a formal series in the ordinary multi-index set; each of its coefficients is a well-defined Hahn sum. Let
\[
 C_\xi=K[[u_1,\ldots,u_n]]/
 (\widehat F_{1,\xi},\ldots,\widehat F_{n,\xi}).
\]
Write $B_\xi$ for the local Artinian factor of $B_K$ at the character $\xi$, and $e_\xi$ for its idempotent.

\begin{theorem}[Local-algebra comparison and exact conservation]\label[theorem]{thm:multiplicity}
For every zero $\xi$ there is a canonical isomorphism
\begin{equation}\label{eq:local-iso}
 B_\xi\simeq C_\xi,
 \qquad e_\xi z_j-\xi_je_\xi\longleftrightarrow u_j.
\end{equation}
In particular all the $C_\xi$ are finite-dimensional, and
\begin{equation}\label{eq:conservation}
 \sum_{\xi:F\sh(\xi)=0}m_\xi=\mu,
 \qquad m_\xi:=\dim_K C_\xi.
\end{equation}
The same dimensions are obtained after scalar extension from $K$ to $\SC$ in the formal local algebras.
\end{theorem}
\begin{proof}
First we prove finiteness of $C_\xi$, rather than presupposing it. The polynomial identity $P_j(z_j)\in(F)$ gives, after Taylor expansion at $\xi$,
\[
 P_j(\xi_j+u_j)=0\quad\hbox{in }C_\xi.
\]
Factor $P_j(T)=(T-\xi_j)^{r_j}Q_j(T)$ with $Q_j(\xi_j)\ne0$. The series $Q_j(\xi_j+u_j)$ is a formal unit, so $u_j^{r_j}=0$ in $C_\xi$. Hence $C_\xi$ is spanned by the finitely many monomials $u^\alpha$ with $\alpha_j<r_j$. It is a local algebra, with residue field $K$.

Taylor expansion defines an algebra map $A_H\to C_\xi$ annihilating the $F_i$ and therefore a map $B_K\to C_\xi$. Modulo the maximal ideal of $C_\xi$, it is evaluation at $\xi$. In a local ring an idempotent is either $0$ or $1$, so the map takes $e_\xi$ to $1$ and the other primary idempotents to $0$. It consequently factors through
\[
 \phi:B_\xi\longrightarrow C_\xi.
\]

For the reverse map, put $n_j=e_\xi z_j-\xi_je_\xi\in B_\xi$. These elements generate the maximal ideal and are jointly nilpotent. Indeed $B_K$ is generated as a $K$-algebra by its coordinates: its chosen basis consists of monomials. Let an ordinary integer $N$ kill the $N$th power of the maximal ideal. Formal substitution $u_j\mapsto n_j$ is then a finite operation on each formal series. By \eqref{eq:finite-taylor}, the natural projection of any $G\in A_H$ into $B_\xi$ is
\[
 \sum_{|\alpha|<N}
 \frac{\partial^\alpha G\sh(\xi)}{\alpha!}n^\alpha.
\]
In particular substitution annihilates every $\widehat F_{i,\xi}$ and induces
\[
 \psi:C_\xi\longrightarrow B_\xi.
\]
The two maps fix the corresponding nilpotent coordinate generators. Both algebras are generated by those elements and $K$, so the maps are inverse. This proves \eqref{eq:local-iso}.

The primary decomposition of a finite-dimensional commutative algebra gives
$B_K\simeq\prod_\xi B_\xi$. Taking dimensions and using \cref{thm:division} yields \eqref{eq:conservation}. Finally the displayed nilpotence relations make every local quotient a quotient of a finite-dimensional polynomial space. Extending its coefficients from $K$ to any field extension commutes with this finite presentation and preserves its dimension. This applies to each set-sized portion needed for $\SC$, and proves the final assertion without forming a set-sized vector space of all class coefficients.
\end{proof}

\begin{corollary}[Simple zeros and a multivariable monad form of Rouch\'e]\label[corollary]{cor:rouche}
Every strictly positive Hahn perturbation of the same ordinary system has total zero multiplicity $\mu$ in the monad of $0$. Its zeros are all simple exactly when
\[
 J_F\sh(\xi)\ne0\quad\hbox{at every zero }\xi.
\]
In that case there are exactly $\mu$ distinct zeros.
\end{corollary}
\begin{proof}
Conservation is \cref{thm:multiplicity}. At a zero, invertibility of the matrix of linear Taylor coefficients implies, by the formal inverse function theorem, that the formal local ideal is the maximal ideal and its quotient has dimension $1$. Conversely, if the quotient has dimension $1$, its defining ideal is the maximal ideal; the $n$ linear parts must span its $n$-dimensional cotangent space, so their determinant is nonzero. Summing the local dimensions proves the last assertion.
\end{proof}
The word Rouch\'e here refers to conservation under a positive Hahn perturbation of the leading system. It is not an unrestricted boundary-norm theorem for arbitrary functions on $\SC^n$.

\section{Spectral invariants without choosing root branches}\label{sec:spectral}

\begin{theorem}[Characteristic polynomials and multiplicity-weighted traces]\label[theorem]{thm:spectral}
For every $G\in A_H$,
\begin{align}
 \det(TI-M_G)
 &=\prod_{\xi:F\sh(\xi)=0}
     \bigl(T-G\sh(\xi)\bigr)^{m_\xi},\label{eq:characteristic}\\
 \Tr(M_G)&=\sum_\xi m_\xi G\sh(\xi),\label{eq:trace}\\
 \det(M_G)&=\prod_\xi G\sh(\xi)^{m_\xi}.\label{eq:norm}
\end{align}
If $S_G=\supp(G)$, the entries of $M_G$ have support in $S_G+M$. For $G\in A_H^+$ they lie in $R$.
\end{theorem}
\begin{proof}
The support assertion follows by applying \cref{thm:division} to each $Gq_a$. In the local factor $B_\xi$, the element $[G]-G\sh(\xi)$ belongs to the nilpotent maximal ideal. Its multiplication operator is therefore nilpotent. Multiplication by $G$ on this factor has sole eigenvalue $G\sh(\xi)$ with algebraic multiplicity $\dim_K B_\xi=m_\xi$. Taking the characteristic polynomial, trace, and determinant of the block decomposition proves the formulas.
\end{proof}

Individual roots can acquire fractional exponents and can exchange under changes of auxiliary descriptions. The matrices above have no such labeling problem. Their entries are produced directly by strong-summation division, and their characteristic coefficients remain invariant under any permutation of the roots. This is particularly useful over high-rank groups, where ordering roots by one real-valued ``size'' would discard information.

\section{Coefficientwise residues and perfect duality}\label{sec:residue}

\subsection{The ordinary residue input and the cycle}
We use the following standard ordinary facts about an isolated holomorphic complete intersection. A sufficiently small isolating neighborhood inside $D$ admits a Grothendieck cycle
\[
 C_f=\{|f_1|=\rho_1,\ldots,|f_n|=\rho_n\},
\]
understood as the local cycle about $0$, with regular positive choices of the radii and orientation determined by
$\dif\arg f_1\wedge\cdots\wedge\dif\arg f_n>0$.
It may be taken inside a neighborhood on which the map germ $f$ is finite. Its ordinary functional
\begin{equation}\label{eq:ordinary-residue}
 \mathcal R_f(g)=\frac1{(2\pi i)^n}
 \int_{C_f}\frac{g\,\dif z_1\wedge\cdots\wedge\dif z_n}{f_1\cdots f_n}
\end{equation}
annihilates $(f)$, and the pairing $(a,b)\mapsto\mathcal R_f(ab)$ on $V$ is perfect. At a simple zero it is evaluation divided by the Jacobian. The same properties apply to positive powers of the $f_i$ in the denominator; in particular a term with one of the denominator divisors absent has zero local residue.

These are classical local-residue and complete-intersection duality facts, not new surcomplex assertions. Algebraic local duality and power-series residues are treated in \cite[Chapters 4--5, 8, and 10]{Kunz}; the contour interpretation, transformation law, and computational consequences are discussed in \cite[Section 1.5]{CattaniDickenstein}. The multivariate residue literature has many additional global results, for example \cite{CattaniDickensteinSturmfels}, which are not being claimed here.

\subsection{Definition on one fixed ordinary cycle}
On a neighborhood of $C_f$, expand
\begin{equation}\label{eq:denom-expansion}
 \frac1{F_i}=\frac1{f_i}\sum_{r\ge0}\left(-\frac{E_i}{f_i}\right)^r.
\end{equation}
No numerical smallness condition is used in this expansion: its correction has strictly positive Hahn support. Define
\begin{equation}\label{eq:H-residue}
 \RF(G)=\frac1{(2\pi i)^n}
 \HInt_{C_f}\frac{G\,\dif z_1\wedge\cdots\wedge\dif z_n}{F_1\cdots F_n}
\end{equation}
by integrating every ordinary coefficient in the product of \eqref{eq:denom-expansion}. Equivalently,
\begin{equation}\label{eq:residue-series}
 \RF(G)=\sum_{r\in\N^n}(-1)^{|r|}
 \frac1{(2\pi i)^n}
 \HInt_{C_f}
 \frac{G E_1^{r_1}\cdots E_n^{r_n}\,
       \dif z_1\wedge\cdots\wedge\dif z_n}
      {f_1^{r_1+1}\cdots f_n^{r_n+1}}.
\end{equation}
The integrals in the coefficients are ordinary complex integrals. The superscript $H$ does not denote a fine-topological Riemann integral over a surcomplex parametrized torus.

\begin{proposition}[Well-defined residue and support]\label[proposition]{prop:residue}
The sum \eqref{eq:residue-series} is strong, is independent of the admissible small local cycle, and defines a $K$-linear functional on $A_H$. It satisfies
\begin{equation}\label{eq:residue-support}
 \supp\RF(G)\subseteq\supp(G)+M,
 \qquad v\RF(G)\ge\vH(G),
 \qquad \RF((F))=0.
\end{equation}
It therefore induces an $R$-linear functional $B_R\to R$ and a $K$-linear functional $B_K\to K$.
\end{proposition}
\begin{proof}
The product expansions have supports in $\supp(G)+M$. For a fixed exponent, \cref{lem:neumann} leaves only finitely many words in the perturbation coefficients, and hence finitely many ordinary meromorphic integrands. Their integrals are complex scalars. This proves strong summability, the support bound, and linearity. Coefficientwise independence of the local cycle is the ordinary residue invariance for those integrands.

For a numerator $G=F_iH$, cancellation of $F_i$ against its strong inverse leaves no denominator divisor $f_i$. Each coefficient is an ordinary holomorphic numerator divided by powers of the other $f_j$, whose local residue on $C_f$ is zero. One can also insert a power of $f_i$ into numerator and denominator and use the ordinary ideal-annihilation property for the system $(f_1^{r_1+1},\ldots,f_n^{r_n+1})$. Thus $\RF(F_iH)=0$. The valuation assertion and descent to the two quotients follow.
\end{proof}

\begin{theorem}[Perfect Hahn residue duality]\label[theorem]{thm:duality}
The map
\begin{equation}\label{eq:duality}
 B_R\longrightarrow\Hom_R(B_R,R),
 \qquad a\longmapsto\bigl(b\mapsto\RF(ab)\bigr)
\end{equation}
is an isomorphism of $B_R$-modules. Thus $B_R$ is a finite free Frobenius $R$-algebra. The analogous pairing over $K$ is perfect, including when the zero set is nonreduced.

In the chosen basis, its matrix
\begin{equation}\label{eq:residue-gram}
 H_F=\bigl(\RF(q_aq_b)\bigr)_{a,b=1}^\mu
\end{equation}
has support in $M$ and determinant in $R^\times$.
\end{theorem}
\begin{proof}
All matrix entries lie in $R$ by \cref{prop:residue}. Their exponent-zero coefficients are exactly $\mathcal R_f(q_aq_b)$, because every occurrence of an $E_i$ has positive support. The ordinary residue pairing is perfect, so the resulting complex matrix $H_0$ has nonzero determinant. Therefore $\det H_F$ has nonzero exponent-zero coefficient and is a unit of $R$. The inverse can either be obtained by the adjugate formula or by a strong geometric inverse about $H_0$; the latter also keeps its entries supported in $M$. The pairing map is represented by $H_F$ and is consequently an isomorphism. Its $B_R$-linearity is the equality $\RF((ca)b)=\RF(a(cb))$.
\end{proof}

\subsection{Normal forms from dual residues}
Let $q^1,\ldots,q^\mu\in B_R$ be the residue-dual basis, so that
\[
 \RF(q_aq^b)=\delta_a^b.
\]
Their coordinates are entries of $H_F^{-1}$ and have support in $M$. Choose their representatives as $R$-linear combinations of the $q_a$.

\begin{corollary}[Contour recovery of the finite algebra]\label[corollary]{cor:residue-reconstruction}
For $G\in A_H$,
\begin{equation}\label{eq:residue-normal-form}
 \NF(G)_a=\RF(Gq^a),
 \qquad (M_G)_{ab}=\RF(Gq_bq^a).
\end{equation}
Thus the residue functional and a fixed ordinary basis recover every normal form and multiplication matrix, without an explicit choice of the Koszul splitting.
\end{corollary}
\begin{proof}
Write the unique remainder of $G$ as $\sum_b c_bq_b$. Pair with $q^a$ and use ideal annihilation. Apply the same argument to $Gq_b$ for the matrix formula.
\end{proof}

\section{The trace--Jacobian and discriminant identities}\label{sec:trace-residue}

\subsection{A finite-parameter comparison principle}
There is a subtlety in transferring an ordinary analytic identity to infinitely many Hahn coefficients: one cannot pretend that an arbitrary Hahn family is a holomorphic family in one ordinary parameter. The correct comparison is coefficientwise and uses finitely many independent parameters at a time.

\begin{lemma}[Finite dependence at one exponent]\label[lemma]{lem:finite-dependence}
Fix an exponent $\nu$. In any expression obtained from the division formulas, multiplication matrices, derivatives, finite determinants, and the residue expansion above, its $t^\nu$ coefficient depends on only finitely many input coefficient functions and finitely many words in their positive perturbation labels. For this one coefficient, the positive labeled terms can be replaced by independent formal variables, and the desired coefficient can be recovered by substituting $s_a=t^{e_a}$ afterwards.
\end{lemma}
\begin{proof}
The formulas consist of finite sums and products and of the geometric series in \eqref{eq:resolvent} and \eqref{eq:denom-expansion}. Every nonconstant iteration inserts at least one exponent in the well-ordered positive set $E$. For the prescribed $\nu$, \cref{lem:neumann} gives finitely many input exponents and finitely many such words. Differentiation, ordinary linear maps, and integration do not change exponents. There are therefore only finitely many relevant labeled holomorphic functions. Assign one independent variable to each of them, keeping labels separate even when their exponents happen to coincide. Ordinary multiplication of the variables records the words. Substituting their positive monomials recovers exactly the finite coefficient sum.
\end{proof}

For completeness, the ordinary analytic comparison needed with those finitely many variables is as follows. Given holomorphic vectors $e_a(z)$, put
\[
 f_s(z)=f(z)+\sum_{a=1}^r s_ae_a(z).
\]
The local analytic algebra
\[
 \C\{z,s\}/(f_s)
\]
is finite over $\C\{s\}$ near $(0,0)$ and is free of rank $\mu$, with basis the $q_a$. Here is the standard justification. Finiteness follows from the analytic finite-mapping theorem because the special fiber is zero-dimensional. The total local ring before quotienting is regular; the $n$ equations have height $n$, so their quotient is Cohen--Macaulay of dimension $r$. Its finite module over the regular parameter ring has full depth and is free. Equivalently, one may use the usual complete-intersection flatness theorem. Nakayama's lemma lifts the chosen basis from the special fiber. These are ordinary finite analytic algebra facts; their analytic foundations are discussed in \cite{Demailly}, and the associated duality algebra in \cite{Kunz}.

On a fixed sufficiently small ordinary residue cycle, the residue integrals and these multiplication matrices are holomorphic in $s$. Their formal Taylor coefficients agree with the degree-filtered versions of the formulas already proved: division by $f_s$ in formal parameter degree gives the same unique remainder in the chosen basis. Restriction to a smaller ordinary polydisc does not alter that remainder, by \cref{sec:division}. Thus this comparison does not require the originally chosen global $h$ to be continuous or the infinite Hahn family to admit ordinary parameter convergence.

\subsection{Trace equals residue against the Jacobian}
\begin{theorem}[Trace--Jacobian identity]\label[theorem]{thm:trace-jacobian}
For every $G\in A_H$,
\begin{equation}\label{eq:trace-jacobian}
 \RF(GJ_F)=\Tr(M_G).
\end{equation}
In particular $\RF(J_F)=\mu$. If all zeros are simple, then
\begin{equation}\label{eq:simple-residue}
 \RF(G)=\sum_{\xi:F\sh(\xi)=0}
 \frac{G\sh(\xi)}{J_F\sh(\xi)}.
\end{equation}
\end{theorem}
\begin{proof}
First recall the ordinary finite-parameter trace identity. At a system with simple zeros, the ordinary local residue formula gives
\[
 \mathcal R_{f_s}(gJ_{f_s})=\sum_{f_s(z)=0}g(z),
\]
while multiplication by $g$ on the reduced finite algebra has those values as its eigenvalues. Hence the right side is its trace. The identity holds also when the fibers are not reduced: adjoin independent small constant shifts to all $n$ equations, use the dense set of regular target values of the finite holomorphic map, and continue the two holomorphic functions of the parameters to the remaining values. Ordinary residues on the fixed isolating cycle vary holomorphically, and so do the multiplication matrices in the finite free basis. This proves the ordinary trace identity with multiplicities, rather than assuming that the initially given parameter family has any simple fiber.

Now fix a Hahn exponent. By \cref{lem:finite-dependence}, only finitely many perturbation labels and numerator coefficient functions affect it in either side of \eqref{eq:trace-jacobian}. Replace these labels by independent ordinary parameters and use the preceding finite analytic family. The ordinary identity holds coefficientwise in their Taylor series. The comparison following \cref{lem:finite-dependence} identifies those coefficients with the formal division and residue expansions. Substituting the positive Hahn monomials therefore proves the desired equality at the chosen exponent. Since the exponent was arbitrary, the Hahn identity follows. Linearity and the same support argument allow an arbitrary Hahn numerator $G$, not just an ordinary one.

Take $G=1$ to obtain $\RF(J_F)=\Tr(I)=\mu$. If all zeros are simple, $J_F$ is a unit in $B_K$, because it is nonzero in every residue field. Apply the identity to the class $[G][J_F]^{-1}$, represented by its finite normal form. Ideal annihilation gives
\[
 \RF(G)=\Tr(M_{G/J_F}).
\]
Finally \cref{thm:spectral} and $m_\xi=1$ give \eqref{eq:simple-residue}.
\end{proof}

\begin{remark}[Multiple zeros]
The perfect pairing and the trace--Jacobian identity do not require simple zeros. The formula with individual terms $G(\xi)/J_F(\xi)$ does. At a multiple zero its denominator is zero, so that expression must not be used. The functional \eqref{eq:H-residue} still exists and retains the nilpotent information that the trace alone cannot recover. No unproved identification of separately localized multiple-zero residues is needed for any result in this article.
\end{remark}

\subsection{Trace degeneracy, the Jacobian norm, and discriminants}
Define the trace Gram matrix
\[
 T_F=\bigl(\Tr M_{q_aq_b}\bigr)_{a,b=1}^\mu,
 \qquad \Delta_F=\det T_F.
\]
This discriminant is relative to the chosen basis; a basis change multiplies it by the square of the change-of-basis determinant.

\begin{theorem}[Discriminant--Jacobian identity]\label[theorem]{thm:discriminant}
In the same basis,
\begin{equation}\label{eq:gram-factorization}
 T_F=H_FM_{J_F},
 \qquad
 \Delta_F=(\det H_F)\det M_{J_F}.
\end{equation}
Consequently
\begin{equation}\label{eq:disc-valuation}
 v(\Delta_F)=v(\det M_{J_F}),
\end{equation}
including the convention $v(0)=+\infty$. The algebra $B_K$ is reduced exactly when $\Delta_F\ne0$, equivalently when all its zeros are simple. In the reduced case,
\begin{equation}\label{eq:disc-sum}
 v(\Delta_F)=\sum_\xi v(J_F\sh(\xi)).
\end{equation}
\end{theorem}
\begin{proof}
The $(a,b)$ entry of $H_FM_{J_F}$ is $\RF(q_aJ_Fq_b)$, which equals $\Tr M_{q_aq_b}$ by \cref{thm:trace-jacobian}. Taking determinants proves \eqref{eq:gram-factorization}; $\det H_F$ is a unit by \cref{thm:duality}, giving \eqref{eq:disc-valuation}.

Over an algebraically closed field of characteristic zero, the trace pairing of a finite commutative algebra is nondegenerate exactly when the algebra is a product of copies of the field. Indeed it is nondegenerate on such a product; a nonzero nilpotent ideal lies in its radical otherwise, since multiplication by a nilpotent element times any other element has trace zero. Thus reducedness is equivalent to $\Delta_F\ne0$. By \cref{thm:multiplicity} and the formal Jacobian criterion in \cref{cor:rouche}, this is also equivalent to all zeros being simple. In that case \eqref{eq:norm} applied to $J_F$ gives the final equality.
\end{proof}

The residue pairing remains unimodular over $R$ even when $\Delta_F$ is infinitesimal or zero. This distinguishes the robust duality pairing from the possibly degenerate trace pairing and explains why the residue basis is preferable near a multiple ordinary zero.

\section{Valuation precision and stability of the finite geometry}\label{sec:precision}
Fix the same $D,f$, and representative basis. Consider two systems
\[
 F=f+E,\qquad \widetilde F=f+\widetilde E
\]
with strictly positive Hahn corrections. Set
\[
 \eta=\min_i\vH(E_i-\widetilde E_i)>0,
\]
with $\eta=+\infty$ when the systems agree. All inequalities below concern the coefficient valuation, not an ordinary analytic norm.

\begin{theorem}[Precision does not deteriorate in the aggregate algebra]\label[theorem]{thm:precision}
For every $G\in A_H$,
\begin{align}
 v\bigl(\NF(G)-\mathcal N_{\widetilde F}(G)\bigr)
 &\ge\vH(G)+\eta,\label{eq:precision-normal}\\
 v\bigl(\RF(G)-\mathcal R_{\widetilde F}(G)\bigr)
 &\ge\vH(G)+\eta.\label{eq:precision-residue}
\end{align}
For a vector or matrix, $v$ means the minimum valuation of its entries. Quotients chosen with the same $h$ satisfy the bound in \eqref{eq:precision-normal} as well. Multiplication matrices obey
\begin{equation}\label{eq:precision-matrix}
 v(M_G^F-M_G^{\widetilde F})\ge\vH(G)+\eta.
\end{equation}
If $G\in A_H^+$, every corresponding characteristic-polynomial coefficient differs by valuation at least $\eta$. The residue Gram matrices, their inverses, the trace Gram matrices, and their determinants also differ by valuation at least $\eta$.
\end{theorem}
\begin{proof}
Use the same ordinary contraction for both systems. Their resolvents satisfy the exact identity
\begin{equation}\label{eq:resolvent-difference}
 S_F-S_{\widetilde F}
 =S_F(\widetilde\delta-\delta)hS_{\widetilde F}.
\end{equation}
Every $S$ preserves the lower support bound, while $\widetilde\delta-\delta$ raises it by at least $\eta$. The maps $p$ and $h$ do not lower it. Applying \eqref{eq:resolvent-difference} to $G$ proves the normal-form and quotient estimates. Applying it to $Gq_a$ gives the matrix estimate.

On the fixed ordinary cycle,
\[
 F_i^{-1}-\widetilde F_i^{-1}
 =\frac{\widetilde F_i-F_i}{F_i\widetilde F_i}.
\]
The inverses on the right have nonnegative Hahn support as meromorphic coefficient data near the cycle. A finite telescoping product of these identities shows that
\[
 v_H\left(\frac G{F_1\cdots F_n}
 -\frac G{\widetilde F_1\cdots\widetilde F_n}\right)
 \ge\vH(G)+\eta.
\]
Coefficientwise integration cannot introduce a smaller exponent, proving \eqref{eq:precision-residue}.

For a nonnegative-support $G$, the entries of both multiplication matrices have nonnegative valuation. Every difference between their characteristic coefficients is a sum of products with at least one changed entry; that entry has valuation at least $\eta$, and all other entries have nonnegative valuation. The same argument applies to the Gram matrices and determinants. For the inverse residue matrices use
\[
 H_F^{-1}-H_{\widetilde F}^{-1}
 =H_F^{-1}(H_{\widetilde F}-H_F)H_{\widetilde F}^{-1},
\]
since both inverse matrices have entries in $R$. For the trace Gram matrices apply \eqref{eq:precision-matrix} to the ordinary numerators $q_aq_b$ and take traces.
\end{proof}

\begin{corollary}[A discriminant threshold for preserving simple zeros]\label[corollary]{cor:disc-stable}
Suppose $F$ has only simple zeros and
\[
 \eta>v(\Delta_F).
\]
Then $\widetilde F$ also has exactly $\mu$ distinct simple zeros, and $\Delta_{\widetilde F}$ has the same leading Hahn term as $\Delta_F$.
\end{corollary}
\begin{proof}
By \cref{thm:precision}, $v(\Delta_{\widetilde F}-\Delta_F)\ge\eta$. The strict inequality prevents cancellation of the leading term of $\Delta_F$, so $\Delta_{\widetilde F}\ne0$. Apply \cref{thm:discriminant} and \cref{thm:multiplicity}.
\end{proof}

\begin{remark}[Root instability versus invariant stability]
The theorem does not assert that a labeling of individual roots is Lipschitz. Roots of a multiple leading system generally involve fractional exponents, and inverse Jacobians can have negative valuation. Normal forms and the total residue nevertheless preserve the input precision because their construction uses only positive-support inverses about ordinary invertible data. The discriminant threshold provides one certified geometric consequence without choosing root branches.
\end{remark}

If the numerator is also changed, the same proof and linearity yield, for either normal forms or total residues, the lower bound
\begin{equation}\label{eq:input-precision}
 \min\bigl\{\vH(G-\widetilde G),\ \vH(G)+\eta\bigr\}
\end{equation}
for the difference between the results for $(F,G)$ and $(\widetilde F,\widetilde G)$. These estimates remain valid at higher-rank exponents where an entire initial segment of the support can be infinite.

\section{Explicit coefficient extraction for monomial leading systems}\label{sec:monomial}
A particularly useful computational case is
\[
 F_i=z_i^{m_i}-A_i(z),\qquad m_i\ge1,
 \qquad\supp A_i\subseteq\Gamma_{>0}.
\]
The ordinary quotient has basis $z^\alpha$ with $0\le\alpha_i<m_i$, and dimension $\mu=\prod_i m_i$.

\begin{corollary}[Strong residue extraction formula]\label[corollary]{cor:coefficient-residue}
For a coherent numerator $G$,
\begin{equation}\label{eq:coefficient-residue}
 \mathcal R_F(G)=
 \sum_{r\in\N^n}
 [z_1^{m_1(r_1+1)-1}\cdots z_n^{m_n(r_n+1)-1}]
 \left(G(z)\prod_{i=1}^n A_i(z)^{r_i}\right).
\end{equation}
Coefficient extraction on the right means ordinary Taylor coefficient extraction on each Hahn coefficient. The whole sum is strongly summable and supported in $\supp(G)+(\bigcup_i\supp A_i)^*$.
\end{corollary}
\begin{proof}
Use an ordinary product of small circles as $C_f$. Expanding
\[
 \frac1{z_i^{m_i}-A_i}=
 \sum_{r_i\ge0}\frac{A_i^{r_i}}{z_i^{m_i(r_i+1)}}
\]
and applying the ordinary one-variable Cauchy coefficient formula in each variable gives \eqref{eq:coefficient-residue}. The support proof is precisely that of \cref{prop:residue}.
\end{proof}
This is the usual residue-extraction mechanism, now justified for an arbitrary positive Hahn perturbation. It is often more explicit than constructing a full Koszul contraction. Together with \eqref{eq:residue-normal-form}, it computes the entire finite algebra.

\section{A coupled quartic over three separated scales}\label{sec:quartic}

\subsection{Elimination and the finite algebra}
Let $a,b,c\in\mm$ and consider
\begin{equation}\label{eq:quartic-system}
 F_1=x^2-y-a,\qquad F_2=y^2-bx-c.
\end{equation}
The leading ordinary system is $(x^2-y,y^2)$, whose only zero is $(0,0)$ and whose quotient is $\C[x]/(x^4)$. Hence $\mu=4$. Eliminating $y$ yields
\begin{equation}\label{eq:quartic}
 P(x)=(x^2-a)^2-bx-c
 =x^4-2ax^2-bx+(a^2-c).
\end{equation}
The main theorem and the spanning monomials show that the full coherent quotient is
\[
 B_K\simeq K[x]/(P),\qquad y=x^2-a,
\]
with basis $1,x,x^2,x^3$. Its multiplication matrices are
\begin{equation}\label{eq:quartic-matrix}
 M_x=\begin{pmatrix}
 0&0&0&c-a^2\\
 1&0&0&b\\
 0&1&0&2a\\
 0&0&1&0
 \end{pmatrix},\qquad M_y=M_x^2-aI.
\end{equation}
The Jacobian is
\begin{equation}\label{eq:quartic-jacobian}
 J_F=4xy-b\equiv4x^3-4ax-b=P'(x).
\end{equation}
In particular the two-variable and eliminated one-variable simple-residue signs agree.

\subsection{The residue pairing and all moments}
In this basis, the residue is the coefficient of $x^3$ in the remainder modulo $P$:
\begin{equation}\label{eq:quartic-residue}
 \RF(g)= [x^3]\operatorname{rem}_{P}\bigl(g(x,x^2-a)\bigr).
\end{equation}
For polynomial $g$, this follows from the ordinary transformation law for elimination, or from the simple-root residue formula and then parameter continuation. For a coherent holomorphic $g$, the right-hand side means its coherent normal form, not polynomial long division applied to an unevaluated infinite expression. The substitution may be made after one restriction to a smaller ordinary $x$-disc on which $x^2$ lies in the ordinary $y$-disc; restriction invariance then applies. The same identity follows coefficientwise from the ordinary residue identities in the independent formal parameters $a,b,c$.

Writing $S_k=\RF(x^k)$ gives
\begin{align}
&S_0=S_1=S_2=0,\qquad S_3=1,\label{eq:moments-init}\\
&S_{k+4}=2aS_{k+2}+bS_{k+1}-(a^2-c)S_k\qquad(k\ge0).\label{eq:moments-recursion}
\end{align}
The first nontrivial values are
\begin{align*}
 S_4&=0,& S_5&=2a,& S_6&=b,\\
 S_7&=3a^2+c,& S_8&=4ab,& S_9&=4a^3+4ac+b^2,\\
 S_{10}&=10a^2b+2bc.&&&&
\end{align*}
Consequently
\begin{equation}\label{eq:quartic-H}
 H_F=\begin{pmatrix}
 0&0&0&1\\
 0&0&1&0\\
 0&1&0&2a\\
 1&0&2a&b
 \end{pmatrix},\qquad\det H_F=1.
\end{equation}
This matrix remains invertible at every degeneration of the root configuration.

The moment generating function is
\begin{equation}\label{eq:moment-generating}
 \sum_{k\ge0}S_k w^k
 =\frac{w^3}{1-2aw^2-bw^3+(a^2-c)w^4}.
\end{equation}
Expanding the geometric denominator proves the finite coefficient identity
\begin{equation}\label{eq:moment-multinomial}
 S_k=\sum_{\substack{r,s,u\ge0\\3+2r+3s+4u=k}}
 \frac{(r+s+u)!}{r!s!u!}(2a)^r b^s(c-a^2)^u.
\end{equation}
For the coherent ordinary exponential numerator $e^x$, expand $1/P$ on a small ordinary $x$-circle as a geometric series about $x^4$. The correction is $2a/x^2+b/x^3+(c-a^2)/x^4$, which has positive Hahn support. Applying the ordinary Cauchy coefficient formula to each term gives the exact strong identity
\begin{equation}\label{eq:quartic-exp}
 \RF(e^x)=\sum_{r,s,u\ge0}
 \frac{(r+s+u)!}{r!s!u!}
 \frac{(2a)^r b^s(c-a^2)^u}{(3+2r+3s+4u)!}.
\end{equation}
For example, grouping the resulting finite coefficient contributions by the degree $k$ of the original numerator gives
\begin{align}\label{eq:quartic-exp-first}
 \RF(e^x)
 &=\frac1{3!}+\frac{2a}{5!}+\frac b{6!}
 +\frac{3a^2+c}{7!}+\frac{4ab}{8!}\notag\\
 &\quad+\frac{4a^3+4ac+b^2}{9!}
 +\frac{10a^2b+2bc}{10!}+\cdots.
\end{align}
This display is ordered by numerator degree, not necessarily by increasing Hahn valuation. The all-order formula \eqref{eq:quartic-exp} is the precise statement. Strong summability follows because every nonconstant factor has positive support; at any exponent only finitely many triples and coefficient choices contribute. In particular
\[
 \st\RF(e^x)=\frac16.
\]

\subsection{Discriminant and genuinely higher-rank root splitting}
Direct calculation gives
\begin{equation}\label{eq:quartic-discriminant}
 \operatorname{disc}(P)
 =-256a^3b^2+256a^2c^2+288ab^2c-27b^4-256c^3.
\end{equation}
Since $\det H_F=1$, \cref{thm:discriminant} identifies this discriminant with both $\det T_F$ and $\det M_{J_F}$.

Now choose
\begin{equation}\label{eq:three-scales}
 a=t,\qquad b=t^\omega,\qquad c=t^{\omega^2},
\end{equation}
in the divisible ordered subgroup generated by $1,\omega,\omega^2$. Then the leading term of \eqref{eq:quartic-discriminant} is $-256a^3b^2$, so all four zeros are simple. This computation compares exponents lexicographically by their $\omega^2$, $\omega$, and finite parts; it is not a single-real-parameter Puiseux approximation.

To display the splitting, write $s=\sqrt a$ and, for $\epsilon=\pm1$, choose
\[
 r_\epsilon^2=b\epsilon s+c.
\]
For each $\delta=\pm1$, a root has leading form
\begin{equation}\label{eq:quartic-root-leading}
 x_{\epsilon,\delta}
 =\epsilon s+\frac{\delta r_\epsilon}{2\epsilon s}(1+o_v(1)),
 \qquad
 y_{\epsilon,\delta}=\delta r_\epsilon(1+o_v(1)).
\end{equation}
Here $o_v(1)$ denotes an element of positive valuation; it is not a sequential asymptotic limit. To justify the two branches near $\epsilon s$, substitute
\[
 x=\epsilon s+\frac{r_\epsilon}{2\epsilon s}W.
\]
The equation becomes
\[
 \left(W+\frac{r_\epsilon}{4a}W^2\right)^2
 -1-\frac b{2\epsilon s r_\epsilon}W=0.
\]
Both displayed correction parameters are infinitesimal in \eqref{eq:three-scales}. At their ordinary value $0$, the equation is $W^2-1=0$, with two simple roots. The ordinary formal implicit-function expansions at $W=\pm1$ evaluate strongly in these positive parameters and give \eqref{eq:quartic-root-leading}. The two choices of $\epsilon$ and two of $\delta$ account for all four roots by \cref{thm:multiplicity}.

At each root,
\[
 \frac{e^{x_{\epsilon,\delta}}}{J_F(x_{\epsilon,\delta},y_{\epsilon,\delta})}
 =\frac1{4\epsilon\delta s r_\epsilon}(1+o_v(1)).
\]
Its valuation is
\[
 -\frac{\omega}{2}-\frac34<0,
\]
so each individual residue is infinitely large. Yet the sum of the four residues has standard part $1/6$, by \eqref{eq:quartic-exp}. The finite coefficientwise contour automatically performs the cancellation across all three scales.

\section{A genuinely transcendental coupled system}\label{sec:transcendental}
Let $a,b\in\mm\setminus\{0\}$ and consider
\begin{equation}\label{eq:trans-system}
 F_1=x^2-ae^y,\qquad F_2=y^2-be^x.
\end{equation}
The exponentials mean the coherent lifts of the ordinary entire functions. All coefficient functions are holomorphic on every fixed ordinary polydisc. The leading system $(x^2,y^2)$ has basis $1,x,y,xy$ and multiplicity $4$. Therefore \cref{thm:multiplicity} gives total multiplicity $4$ in the monad of $(0,0)$, without requiring algebraic elimination of either exponential.

\subsection{Four simple roots}
At a zero, $x$ and $y$ are nonzero because $a,b$ and the two exponential values are nonzero. Moreover
\begin{equation}\label{eq:trans-jacobian}
 J_F=4xy-ab e^{x+y}=xy(4-xy),
\end{equation}
since $x^2y^2=ab e^{x+y}$. The product $xy$ is infinitesimal, so $4-xy$ is a unit. Thus every zero is simple, and there are exactly four.

For explicit leading branches, choose $s=\sqrt a$, $r=\sqrt b$ and set $x=sU$, $y=rV$. The equations become
\[
 U^2=e^{rV},\qquad V^2=e^{sU}.
\]
At $(s,r)=(0,0)$ there are four simple ordinary solutions $(U,V)=(\epsilon,\delta)$ with $\epsilon,\delta\in\{\pm1\}$. Ordinary analytic implicit-function expansions in the two independent parameters $s,r$ evaluate strongly at our Hahn values. They give
\begin{equation}\label{eq:trans-leading}
 x_{\epsilon,\delta}=\epsilon\sqrt a(1+o_v(1)),
 \qquad
 y_{\epsilon,\delta}=\delta\sqrt b(1+o_v(1)).
\end{equation}
This works equally well for $a=t$, $b=t^\omega$ or for positive supports with infinitely many exponents.

\subsection{An exact two-scale transcendental residue}
Apply \cref{cor:coefficient-residue} with numerator $G=e^{x+y}$. The denominator expansion gives the coefficient of $a^r b^s$ as the ordinary double residue of
\[
 \frac{e^{(s+1)x+(r+1)y}}{x^{2r+2}y^{2s+2}}\,\dif x\wedge\dif y.
\]
Cauchy coefficient extraction yields the exact formula
\begin{equation}\label{eq:trans-residue}
 \boxed{\displaystyle
 \RF(e^{x+y})=
 \sum_{r,s\ge0}a^r b^s
 \frac{(s+1)^{2r+1}(r+1)^{2s+1}}
 {(2r+1)!(2s+1)!}.}
\end{equation}
It is strongly summable for every pair of nonzero infinitesimal $a,b$. Its first terms, grouped by total degree in these two parameters, are
\begin{equation}\label{eq:trans-residue-first}
 1+\frac{a+b}{3}+\frac{a^2+b^2}{40}+\frac{16ab}{9}
 +\text{terms of total parameter degree at least }3.
\end{equation}
Parameter degree is again a bookkeeping convention, not an assertion that the displayed terms are already in Hahn order.

By \cref{thm:trace-jacobian}, the same quantity is the finite sum
\begin{equation}\label{eq:trans-root-sum}
 \sum_{\epsilon,\delta=\pm1}
 \frac{e^{x_{\epsilon,\delta}+y_{\epsilon,\delta}}}
 {x_{\epsilon,\delta}y_{\epsilon,\delta}
  (4-x_{\epsilon,\delta}y_{\epsilon,\delta})}.
\end{equation}
Each summand has leading value
\[
 \frac{\epsilon\delta}{4\sqrt{ab}},
\]
which is infinitely large. The exact total has standard part $1$. This example shows why the residue functional cannot be replaced by a procedure that discards all infinite individual terms. Their cancellation is part of the mathematical object, not an error to be removed.

\subsection{A nonreduced degeneration}
The division and perfect-pairing theorems also apply when $a=0$ or $b=0$. For instance, if $a=0$ and $b\ne0$, the zeros are $(0,\pm\sqrt b)$, each of multiplicity $2$. To see the multiplicity directly, the second equation has nonzero $y$-derivative at each such point, so it solves formally for $y$ in terms of $x$; the remaining equation is $x^2=0$. The total multiplicity is still $4$, and \eqref{eq:trans-residue} still defines the residue after setting $a=0$. The simple-root expression \eqref{eq:trans-root-sum} must not be used at this degeneration, since its Jacobians vanish.

\section{Scope, novelty assessment, and remaining questions}\label{sec:status}

\subsection{The part of the manuscript's program resolved here}
Within the common-domain Hahn-coherent category of \cite{Companion}, the following precise problem has been answered: an arbitrary isolated ordinary complete intersection admits a full positive-Hahn deformation theory that keeps division on one fixed coefficient domain and preserves its finite algebra and total multiplicity. The proof works in every finite number of variables. It is not restricted to triangular systems, polynomial perturbations, finitely generated support monoids, or simple ordinary roots.

The principal extension statements are \cref{thm:division,thm:points,thm:multiplicity,thm:duality,thm:trace-jacobian,thm:precision}. Their combination provides a finite-mapping theorem in the algebraic sense, an exact geometric multiplicity comparison, and residue and discriminant computations stable under positive-valuation errors. It resolves the isolated, zero-dimensional complete-intersection portion of the proposed several-variable theory. It does not purport to settle all positive-dimensional analytic geometry mentioned in the source.

\subsection{What the literature check supports}
The check for this draft included the attached text, Alling's published framework \cite{Alling}, Hahn-field algebra and closedness \cite{Neumann,Poonen}, the general analytic-structure theory \cite{CluckersLipshitz}, the perturbation lemma \cite{Crainic}, and ordinary local residue and finite-algebra theory \cite{Demailly,Kunz,CattaniDickenstein}. The published surreal integration work of Costin and Ehrlich \cite{CostinEhrlich} was also checked to avoid representing broad extension and integration questions as untouched. Their results concern a different extension and integration program, including substantial resurgent classes; our coefficientwise local contour construction does not replace or subsume it.

The exact package consisting of full $\OO(D)((t^\Gamma))$ coefficients, one unchanged ordinary domain, the explicit $S+E^*$ support bound, comparison with all actual surcomplex zeros, and the residue/precision consequences was not located in the checked material. That supports presenting it as a candidate new theorem package in this particular framework. It does \emph{not} establish that no equivalent theorem can be extracted from the broader literature. In particular, a general analytic-structure or deformation-theory result could have overlapping consequences under translated hypotheses. An exhaustive priority comparison is separate from checking the proofs given here.

The named open target comes from the supplied unpublished manuscript. We have not identified a published conjecture whose established name and exact hypotheses coincide with \cref{thm:division} or \cref{thm:multiplicity}. Accordingly this article is a construction of substantial extension theorems, not an announcement that a classical conjecture has been solved.

\subsection{Boundaries of the proofs}
Several restrictions do mathematical work. The leading zero set must be isolated, so that its ordinary quotient is finite-dimensional. There must be exactly a complete-intersection system; arbitrary redundant generators cannot simply be fed into the same positive-degree Koszul exactness argument. The perturbations must have positive support and a common ordinary domain. Negative-support corrections can change the leading system and its multiplicity, and unrelated coefficient germs need not retain one fixed domain under iteration.

The contour functional is coefficientwise. Nothing here manufactures a fine-topological Riemann integral or uses sequential compactness on $\SC$. All infinite sums are justified by finite contributions at each exponent. The precision results concern Hahn coefficients and finite-algebra invariants; they are not estimates in an unmentioned real-valued norm.

\subsection{Four concrete next questions}
\paragraph{Positive-dimensional complete intersections.}
When the ordinary quotient has infinite dimension, the perturbation lemma still has formal meaning if a suitable contraction exists, but the finite spectral and residue arguments above no longer apply. A useful next target is a coherent module over an ordinary parameter domain, with a locally free finite-rank direct image and transition maps that have globally compatible support. Merely applying the present theorem at each ordinary point would not prove such compatibility.

\paragraph{General finite analytic ideals.}
An isolated ideal with more than $n$ specified generators has syzygies not encoded by the complete-intersection Koszul complex. A finite free resolution, together with a perturbation of its full differential satisfying the square-zero equations, is a plausible replacement. Arbitrarily perturbing generators without simultaneously controlling these syzygies need not preserve length, so the correct compatibility hypothesis is part of the problem.

\paragraph{Effective division without arbitrary splittings.}
For monomial leading systems, \cref{cor:coefficient-residue} is explicit. For general holomorphic systems, one can seek residue or integral operators with quantitatively controlled ordinary complexity. Our theorem proves existence and support control; it does not turn an arbitrary presentation of a holomorphic function into an executable finite procedure.

\paragraph{Sectorial and transcendental globalizations.}
The present finite geometry is local over an ordinary monad. It does not solve the source's sectorial summation and overlap questions or define a new global surcomplex exponential. Any globalization must specify which coefficient descriptions are identified across domains and prove compatibility of its contour functionals. The rigid common-domain construction here supplies a local comparison object, not that global compatibility theorem.

\section{Conclusion}
A multiple ordinary zero does not merely leave behind a list of displaced surcomplex points. Under a Hahn-coherent perturbation it leaves behind a finite free algebra of unchanged rank, and that algebra is the correct carrier of multiplicity, residues, and spectral information. The decisive construction is to extend one ordinary Koszul contraction coefficientwise and sum only along a positive, well-ordered perturbation support. This keeps all division coefficients on the original domain while avoiding unjustified convergence assumptions.

The geometric comparison proves that the finite algebra has exactly the intended surcomplex points and exactly their formal local multiplicities. Ordinary residue duality then survives over the possibly non-Noetherian Hahn valuation ring, while the trace--Jacobian identity gives a genuine multidimensional counting formula. Finally, precision is preserved for these aggregate invariants even when individual roots split at fractional or widely separated scales and individual residues become infinite. These results supply a substantial zero-dimensional several-variable continuation of the attached framework.

\appendix
\section{Dependency and hypothesis audit}\label{app:audit}
The proof uses a small number of established inputs. The extension steps and their dependencies can be read in the following order.
\begin{center}\small
\begin{tabular}{L{0.24\textwidth}L{0.32\textwidth}L{0.34\textwidth}}
\toprule
Stage & Established input & Extension step proved here\\
\midrule
Coefficient calculus & Positive-support Neumann lemma & Strong evaluation, global divided differences, and support-preserving linear maps.\\
Ordinary finite algebra & Regular local rings, Koszul exactness, Cartan B & One global algebraic contraction on $\OO(D)$.\\
Perturbed division & Homological perturbation identities & Strong resolvent on full Hahn modules; fixed-domain quotients with support $S+E^*$.\\
Geometry & Hahn algebraic closedness; finite Artinian algebras & Characters are actual evaluations; every local factor equals the formal germ algebra at a displaced zero.\\
Residues & Ordinary Grothendieck duality and local cycles & Perfect pairing over $R$, finite-parameter trace comparison, discriminant and valuation estimates.\\
\bottomrule
\end{tabular}
\end{center}

The assumptions can also be audited by the step where each enters. Isolatedness gives finite $V$ and nilpotent coordinate classes. The complete-intersection condition gives positive-degree Koszul exactness. The common ordinary domain permits one global $h$. Positive supports give the geometric inverse and prohibit support loss. Divisibility of $\Gamma$ and algebraic closedness of $\C$ ensure that all finite-algebra characters have coordinates in the chosen $K$. Ordinary residue duality gives the nonzero leading determinant of $H_F$.

None of these steps uses a discrete valuation, an Archimedean value group, a Noetherian Hahn valuation ring, eventual convergence of finite partial sums in the fine topology, or a choice of individual root branches.

\section{Reproducible finite checks}\label{app:checks}
The archive contains \texttt{verify\_examples.py} and its actual output \texttt{verification\_report.txt}. The script uses exact symbolic arithmetic in SymPy. It verifies elimination and both matrix equations for \eqref{eq:quartic-system}; the companion characteristic polynomial; the equality of the two Jacobians; the residue Gram matrix and its determinant; thirteen trace--Jacobian identities; the multinomial moment formula through degree $15$; the discriminant factorization; the leading discriminant exponent in the rank-three example; and twenty-five independently extracted coefficients in \eqref{eq:trans-residue}.

The computation was run with Python 3.13.5 and SymPy 1.14.0. Every assertion passed. These are finite consistency checks of the worked formulas. They are not a machine proof of strong summability, ordinary Cartan vanishing, the abstract deformation theorem, or the local-algebra comparison. Those are justified by the mathematical arguments and identified classical inputs in the body of the article.

The \LaTeX{} source is self-contained, including its bibliography. Compile it twice with \texttt{pdflatex} to resolve references and the table of contents. No external images, font files, or bibliography database are required.

\begin{thebibliography}{99}

\bibitem{Companion}
\emph{Surcomplex Analysis: Infinitesimal Calculus, Hahn-Coherent Holomorphy, and Contour Theory}.
Unpublished manuscript supplied with the request, file
\texttt{surcomplex\_analysis(3).tex}, dated September 21, 2026.
The definitions of Hahn-coherent data and the final several-variable research direction are the starting point of this article.

\bibitem{Alling}
N.~L. Alling,
\emph{Foundations of Analysis over Surreal Number Fields},
North-Holland Mathematics Studies, vol.~141, North-Holland, 1987.
Publisher record: \url{https://shop.elsevier.com/books/foundations-of-analysis-over-surreal-number-fields/alling/978-0-444-70226-5}.

\bibitem{Neumann}
B.~H. Neumann,
On ordered division rings,
\emph{Transactions of the American Mathematical Society} \textbf{66} (1949), 202--252.
\href{https://doi.org/10.1090/S0002-9947-1949-0032593-5}{doi:10.1090/S0002-9947-1949-0032593-5}.

\bibitem{Poonen}
B. Poonen,
Maximally complete fields,
\emph{L'Enseignement Math\'ematique} \textbf{39} (1993), 87--106.
Author's text: \url{https://math.mit.edu/~poonen/papers/amsval.pdf}.

\bibitem{CluckersLipshitz}
R. Cluckers and L. Lipshitz,
Fields with analytic structure,
\emph{Journal of the European Mathematical Society} \textbf{13} (2011), no.~4, 1147--1223.
\href{https://doi.org/10.4171/JEMS/278}{doi:10.4171/JEMS/278}.
Preprint: \href{https://arxiv.org/abs/0908.2376}{arXiv:0908.2376}.

\bibitem{Crainic}
M. Crainic,
\emph{On the perturbation lemma, and deformations},
preprint, 2004.
\href{https://arxiv.org/abs/math/0403266}{arXiv:math/0403266}.
The sign of the contracting homotopy in the present article is chosen so that $dh+hd=1-ip$.

\bibitem{Demailly}
J.-P. Demailly,
\emph{Complex Analytic and Differential Geometry},
online monograph, version of June 21, 2012.
\url{https://www-fourier.univ-grenoble-alpes.fr/~demailly/manuscripts/agbook.pdf}.

\bibitem{Kunz}
E. Kunz,
\emph{Residues and Duality for Projective Algebraic Varieties},
University Lecture Series, vol.~47, American Mathematical Society, 2008.
\href{https://doi.org/10.1090/ulect/047}{doi:10.1090/ulect/047}.

\bibitem{CattaniDickenstein}
E. Cattani and A. Dickenstein,
Introduction to residues and resultants,
in A. Dickenstein and I.~Z. Emiris (eds.),
\emph{Solving Polynomial Equations: Foundations, Algorithms, and Applications},
Algorithms and Computation in Mathematics, vol.~14, Springer, 2005, pp.~1--61.
\href{https://doi.org/10.1007/3-540-27357-3_1}{doi:10.1007/3-540-27357-3\_1}.
Author-hosted book contents: \url{https://mate.dm.uba.ar/~alidick/contents.html}.

\bibitem{CattaniDickensteinSturmfels}
E. Cattani, A. Dickenstein, and B. Sturmfels,
Residues and resultants,
\emph{Journal of Mathematical Sciences, the University of Tokyo} \textbf{5} (1998), 119--148.
\href{https://arxiv.org/abs/alg-geom/9702001}{arXiv:alg-geom/9702001}.
Journal record: \url{https://www.ms.u-tokyo.ac.jp/journal/abstract_e/jms050106_e.html}.

\bibitem{CostinEhrlich}
O. Costin and P. Ehrlich,
Integration on the surreals,
\emph{Advances in Mathematics} \textbf{452} (2024), article 109823.
\href{https://doi.org/10.1016/j.aim.2024.109823}{doi:10.1016/j.aim.2024.109823}.
Preprint: \href{https://arxiv.org/abs/2208.14331}{arXiv:2208.14331}.

\end{thebibliography}
\end{document}
